\documentclass[11pt]{article}

\usepackage{amsmath, amssymb, amsthm}
\usepackage{geometry}
\usepackage{graphicx}
\usepackage{enumitem}

\geometry{margin=1in}

\title{CSE 400: Stochastic Processes}
\author{Lecture 26: SSS, WSS, and Ergodic Processes \\ Scribe Notes}
\date{April 2026}

\begin{document}

\maketitle

\section{Introduction}

A \textbf{stochastic process} is a collection of random variables indexed by time:
\[
X(t), \quad t \in T
\]

Each outcome produces a \textbf{realization} (sample function).

\subsection{Key Concepts}
\begin{itemize}
    \item \textbf{Random Variable:} $X(t)$ at fixed $t$
    \item \textbf{Realization:} Function of time for fixed outcome
    \item \textbf{Ensemble:} Collection of all realizations
\end{itemize}

\section{Characterization of Stochastic Process}

\subsection{Distribution Functions}

\begin{itemize}
    \item First-order PDF: $f_X(x,t)$
    \item Second-order PDF: $f_X(x_1, t_1; x_2, t_2)$
    \item General $n$-th order PDF
\end{itemize}

\section{Statistical Moments}

\subsection{Mean Function}

\[
\mu_X(t) = E[X(t)]
\]

\subsection{Autocorrelation Function}

\[
R_X(t_1, t_2) = E[X(t_1)X(t_2)]
\]

\subsection{Covariance Function}

\[
C_X(t_1, t_2) = E[(X(t_1) - \mu_X(t_1))(X(t_2) - \mu_X(t_2))]
\]

Relation:
\[
C_X(t_1,t_2) = R_X(t_1,t_2) - \mu_X(t_1)\mu_X(t_2)
\]

\section{Properties of Autocorrelation}

\begin{itemize}
    \item Symmetry:
    \[
    R_X(t_1,t_2) = R_X(t_2,t_1)
    \]
    
    \item Maximum at origin:
    \[
    R_X(t,t) \geq R_X(t_1,t_2)
    \]
    
    \item Non-negative:
    \[
    R_X(t,t) \geq 0
    \]
\end{itemize}

\section{Positive Definite Function}

A function $R_X$ is \textbf{positive definite} if:
\[
\sum_{i=1}^{n} \sum_{j=1}^{n} a_i a_j R_X(t_i, t_j) \geq 0
\]

\subsection{Matrix Representation}

\[
\mathbf{R} =
\begin{bmatrix}
R(t_1,t_1) & \cdots & R(t_1,t_n) \\
\vdots & \ddots & \vdots \\
R(t_n,t_1) & \cdots & R(t_n,t_n)
\end{bmatrix}
\]

\subsection{Condition}
\[
\mathbf{a}^T \mathbf{R} \mathbf{a} \geq 0
\]

\section{Stationary Processes}

\subsection{Strict Sense Stationary (SSS)}

A process is SSS if all statistical properties are invariant under time shift:
\[
f_X(x_1, t_1; \dots; x_n, t_n) = f_X(x_1, t_1+\tau; \dots; x_n, t_n+\tau)
\]

\subsection{Wide Sense Stationary (WSS)}

A process is WSS if:

\begin{itemize}
    \item Mean is constant:
    \[
    \mu_X(t) = \mu
    \]
    
    \item Autocorrelation depends only on time difference:
    \[
    R_X(t_1,t_2) = R_X(\tau), \quad \tau = t_2 - t_1
    \]
\end{itemize}

\subsection{SSS vs WSS}

\begin{itemize}
    \item SSS $\Rightarrow$ WSS
    \item WSS $\nRightarrow$ SSS (in general)
\end{itemize}

\subsection{Gaussian Process Special Case}

For Gaussian processes:
\[
\text{WSS} \Rightarrow \text{SSS}
\]

\section{Example: White Noise Process}

\begin{itemize}
    \item Mean:
    \[
    \mu = 0
    \]
    
    \item Autocorrelation:
    \[
    R_X(\tau) = \sigma^2 \delta(\tau)
    \]
\end{itemize}

\section{Ergodic Process}

A process is \textbf{ergodic} if time averages equal ensemble averages.

\subsection{Time Average}

\[
\overline{X} = \lim_{T \to \infty} \frac{1}{T} \int_{-T/2}^{T/2} X(t) dt
\]

\subsection{Ergodicity Condition}

\[
\overline{X} = E[X(t)]
\]

\section{Example: Mean and Autocorrelation Derivation}

\subsection{Mean}
\[
\mu_X(t) = E[X(t)]
\]

\subsection{Autocorrelation}
\[
R_X(t_1,t_2) = E[X(t_1)X(t_2)]
\]

Steps:
\begin{enumerate}
    \item Express process
    \item Compute expectation
    \item Use independence (if applicable)
\end{enumerate}

\section{Important Mathematical Tools}

\begin{itemize}
    \item Expectation operator linearity
    \item Independence:
    \[
    E[XY] = E[X]E[Y]
    \]
    \item Delta function properties
\end{itemize}

\section{Exam-Oriented Summary}

\subsection{Key Definitions}

\begin{itemize}
    \item Stochastic Process = Collection of RVs
    \item Mean: $E[X(t)]$
    \item Autocorrelation: $E[X(t_1)X(t_2)]$
    \item Covariance: $R - \mu^2$
\end{itemize}

\subsection{Stationarity}

\begin{itemize}
    \item SSS: Full distribution invariant
    \item WSS: Mean constant + autocorrelation depends only on $\tau$
\end{itemize}

\subsection{Important Results}

\begin{itemize}
    \item SSS $\Rightarrow$ WSS
    \item Gaussian: WSS $\Rightarrow$ SSS
\end{itemize}

\subsection{Autocorrelation Properties}

\begin{itemize}
    \item Symmetric
    \item Positive definite
    \item Maximum at origin
\end{itemize}

\subsection{Ergodicity}

\begin{itemize}
    \item Time average = Ensemble average
\end{itemize}

\subsection{Must-Know Examples}

\begin{itemize}
    \item White noise
    \item Mean and autocorrelation derivation
\end{itemize}

\section{Final Remarks}

This lecture builds the foundation for:
\begin{itemize}
    \item Signal processing
    \item Random signal analysis
    \item Communication systems
\end{itemize}

\end{document}